\documentclass[12pt]{article}
\usepackage[utf8]{inputenc}
\usepackage[T1]{fontenc}
\usepackage{lmodern}
\usepackage{amsmath,amssymb,amsthm}
\usepackage{graphicx}
\usepackage[margin=1in]{geometry}
\usepackage{xcolor}
\usepackage{hyperref}
\usepackage{booktabs}
\usepackage{natbib}
\usepackage{lipsum}
\usepackage{float}
\usepackage{caption}
\usepackage{subcaption}
\usepackage{enumitem}
\usepackage{microtype}

\hypersetup{
  colorlinks=true,
  linkcolor=blue!70!black,
  citecolor=green!50!black,
  urlcolor=blue!80!black
}

\newtheorem{theorem}{Theorem}[section]
\newtheorem{lemma}[theorem]{Lemma}
\newtheorem{definition}[theorem]{Definition}

\title{\textbf{Adaptive Neural Architecture Search\\via Gradient-Based Optimization}}
\author{
  Alice M.\ Johnson\textsuperscript{1} \and
  Robert K.\ Chen\textsuperscript{2} \and
  Maria L.\ Santos\textsuperscript{1}\\
  \textsuperscript{1}Department of Computer Science, Stanford University\\
  \textsuperscript{2}Google DeepMind\\
  \texttt{\{ajohnson, msantos\}@cs.stanford.edu, rchen@deepmind.com}
}
\date{\today}

\begin{document}

\maketitle

\begin{abstract}
Neural Architecture Search (NAS) has emerged as a powerful paradigm for automating the design of deep neural networks. However, existing approaches often suffer from prohibitive computational costs and limited generalization across tasks. In this paper, we propose \textit{AdaptiveNAS}, a novel gradient-based framework that leverages differentiable architecture parameters combined with a hierarchical search space to efficiently discover high-performing architectures. Our method introduces a multi-objective optimization criterion that simultaneously considers accuracy, latency, and model size. Extensive experiments on CIFAR-10, CIFAR-100, and ImageNet demonstrate that AdaptiveNAS achieves state-of-the-art performance while reducing search costs by $4.7\times$ compared to previous methods. We further show that architectures discovered by our method transfer effectively to downstream tasks including object detection and semantic segmentation.
\end{abstract}

\section{Introduction}

Deep neural networks have achieved remarkable success across a wide range of applications, from image classification~\citep{he2016deep} to natural language processing~\citep{vaswani2017attention}. However, the design of neural network architectures remains a labor-intensive process that requires significant domain expertise. Neural Architecture Search (NAS) aims to automate this process by searching over a space of possible architectures to find those that maximize performance on a given task.

Early NAS methods relied on reinforcement learning~\citep{zoph2017neural} or evolutionary algorithms~\citep{real2019regularized} to explore the architecture space. While these approaches demonstrated the potential of automated architecture design, they required thousands of GPU hours to complete a single search, making them impractical for many applications. Recent work has focused on reducing the computational cost of NAS through weight sharing~\citep{pham2018efficient} and differentiable relaxation~\citep{liu2019darts}.

In this paper, we propose AdaptiveNAS, a novel gradient-based architecture search framework that addresses two key limitations of existing methods:

\begin{enumerate}[label=(\roman*)]
  \item \textbf{Search efficiency}: Our method uses a hierarchical search space that decomposes the architecture design problem into manageable sub-problems, enabling efficient gradient-based optimization.
  \item \textbf{Multi-objective optimization}: We introduce a differentiable surrogate for hardware-aware metrics, allowing simultaneous optimization of accuracy, latency, and model size.
\end{enumerate}

Our main contributions are summarized as follows:
\begin{itemize}
  \item We propose a hierarchical differentiable search space that captures both cell-level and network-level architectural decisions.
  \item We develop a novel multi-objective optimization strategy that balances accuracy with computational efficiency.
  \item We achieve state-of-the-art results on CIFAR-10 (97.8\%), CIFAR-100 (85.3\%), and ImageNet (79.2\% top-1) with significantly reduced search costs.
\end{itemize}

\section{Related Work}

\subsection{Neural Architecture Search}

The field of NAS has evolved rapidly since the seminal work of Zoph and Le~\citep{zoph2017neural}, who used a recurrent neural network controller trained with reinforcement learning to generate architecture descriptions. Subsequent work by Real et al.~\citep{real2019regularized} demonstrated that evolutionary methods could achieve comparable results with improved stability.

\subsection{Differentiable Architecture Search}

DARTS~\citep{liu2019darts} introduced a continuous relaxation of the discrete architecture space, enabling gradient-based optimization. This approach reduced search costs from thousands of GPU hours to a single GPU day. However, DARTS is known to suffer from performance collapse, where the search converges to architectures dominated by skip connections.

\subsection{Hardware-Aware NAS}

Several recent works have incorporated hardware constraints directly into the search process. MnasNet~\citep{tan2019mnasnet} used a multi-objective reward function that balances accuracy and latency. FBNet~\citep{wu2019fbnet} extended differentiable search to include latency-aware optimization through a lookup table approach.

\section{Method}

\subsection{Problem Formulation}

Let $\mathcal{A}$ denote the space of all possible architectures. Our goal is to find an architecture $a^* \in \mathcal{A}$ that minimizes a multi-objective loss function:

\begin{equation}
  a^* = \arg\min_{a \in \mathcal{A}} \; \mathcal{L}_{\text{val}}(w^*(a), a) + \lambda_1 \cdot \text{LAT}(a) + \lambda_2 \cdot \text{SIZE}(a)
\end{equation}

where $w^*(a) = \arg\min_w \mathcal{L}_{\text{train}}(w, a)$ are the optimal weights for architecture $a$, $\text{LAT}(a)$ is the inference latency, and $\text{SIZE}(a)$ is the number of parameters.

\subsection{Hierarchical Search Space}

We decompose the search space into two levels. At the \textit{cell level}, we search for computational cells that serve as building blocks. At the \textit{network level}, we determine the arrangement and connectivity of these cells.

\begin{definition}[Computational Cell]
A computational cell is a directed acyclic graph (DAG) $G = (V, E)$ where each node $v_i \in V$ represents a latent feature map and each edge $(i, j) \in E$ is associated with an operation $o_{ij} \in \mathcal{O}$.
\end{definition}

The set of candidate operations $\mathcal{O}$ includes:
\begin{itemize}
  \item $3 \times 3$ and $5 \times 5$ separable convolutions
  \item $3 \times 3$ and $5 \times 5$ dilated separable convolutions
  \item $3 \times 3$ max pooling and average pooling
  \item Identity (skip connection)
  \item Zero (no connection)
\end{itemize}

\subsection{Differentiable Relaxation}

Following DARTS, we relax the categorical choice of operations to a continuous distribution using a softmax over architecture parameters $\alpha$:

\begin{equation}
  \bar{o}_{ij}(x) = \sum_{o \in \mathcal{O}} \frac{\exp(\alpha_o^{ij})}{\sum_{o' \in \mathcal{O}} \exp(\alpha_{o'}^{ij})} \cdot o(x)
\end{equation}

\begin{theorem}[Convergence Guarantee]
Under mild regularity conditions on $\mathcal{L}_{\text{val}}$ and $\mathcal{L}_{\text{train}}$, the bilevel optimization converges to a stationary point at a rate of $O(1/\sqrt{T})$ where $T$ is the number of iterations.
\end{theorem}

\section{Experiments}

\subsection{Experimental Setup}

We evaluate AdaptiveNAS on three benchmark datasets: CIFAR-10, CIFAR-100, and ImageNet. For the search phase, we use a proxy task on CIFAR-10 with 8 cells and 16 initial channels. The search is conducted for 50 epochs using SGD with momentum 0.9 and weight decay $3 \times 10^{-4}$.

\subsection{Results on CIFAR-10 and CIFAR-100}

Table~\ref{tab:cifar_results} summarizes the comparison with state-of-the-art methods on CIFAR-10 and CIFAR-100.

\begin{table}[H]
\centering
\caption{Comparison of architectures on CIFAR-10 and CIFAR-100.}
\label{tab:cifar_results}
\begin{tabular}{@{}lccccc@{}}
\toprule
\textbf{Method} & \textbf{CIFAR-10} & \textbf{CIFAR-100} & \textbf{Params} & \textbf{Search Cost} & \textbf{Search}\\
 & \textbf{Error (\%)} & \textbf{Error (\%)} & \textbf{(M)} & \textbf{(GPU days)} & \textbf{Method} \\
\midrule
NASNet-A        & 2.65 & 17.81 & 3.3 & 1800 & RL \\
AmoebaNet-A     & 2.55 & 18.93 & 3.2 & 3150 & Evolution \\
ENAS            & 2.89 & 19.43 & 4.6 & 0.5  & RL \\
DARTS (2nd)     & 2.76 & 17.54 & 3.3 & 1.0  & Gradient \\
PC-DARTS        & 2.57 & 17.11 & 3.6 & 0.1  & Gradient \\
\midrule
\textbf{AdaptiveNAS} & \textbf{2.21} & \textbf{14.72} & \textbf{3.4} & \textbf{0.2} & \textbf{Gradient} \\
\bottomrule
\end{tabular}
\end{table}

\subsection{Results on ImageNet}

We transfer the best cell architecture discovered on CIFAR-10 to ImageNet by constructing a larger network with 14 cells and 48 initial channels. As shown in Table~\ref{tab:imagenet_results}, AdaptiveNAS achieves competitive top-1 accuracy while maintaining reasonable model complexity.

\begin{table}[H]
\centering
\caption{ImageNet classification results.}
\label{tab:imagenet_results}
\begin{tabular}{@{}lcccc@{}}
\toprule
\textbf{Method} & \textbf{Top-1 (\%)} & \textbf{Top-5 (\%)} & \textbf{Params (M)} & \textbf{FLOPs (M)} \\
\midrule
MobileNetV2    & 72.0 & 91.0 & 3.4 & 300 \\
NASNet-A       & 74.0 & 91.6 & 5.3 & 564 \\
DARTS          & 73.3 & 91.3 & 4.7 & 574 \\
ProxylessNAS   & 75.1 & 92.5 & 7.1 & 465 \\
\midrule
\textbf{AdaptiveNAS} & \textbf{79.2} & \textbf{94.5} & \textbf{5.8} & \textbf{490} \\
\bottomrule
\end{tabular}
\end{table}

\section{Analysis}

\subsection{Ablation Study}

To understand the contribution of each component, we conduct ablation experiments on CIFAR-10. Removing the hierarchical search space increases the error rate from 2.21\% to 2.58\%. Removing the multi-objective loss leads to architectures with 2.3$\times$ more parameters while providing only marginal accuracy improvements. These results confirm that both components are essential for achieving the best accuracy-efficiency trade-off.

\subsection{Search Cost Analysis}

Our method completes the architecture search in approximately 4.8 GPU hours on a single NVIDIA V100, which is $4.7\times$ faster than DARTS and orders of magnitude faster than RL-based or evolutionary methods. The efficiency gain primarily comes from the hierarchical decomposition, which reduces the number of architecture parameters by 60\%.

\section{Conclusion}

We presented AdaptiveNAS, a gradient-based neural architecture search framework that achieves state-of-the-art results across multiple benchmarks while significantly reducing search costs. Our hierarchical search space and multi-objective optimization strategy enable the discovery of architectures that are both accurate and efficient. Future work will explore extending AdaptiveNAS to additional domains such as natural language processing and speech recognition, as well as incorporating additional hardware targets.

\section*{Acknowledgments}

This work was supported in part by NSF Grant IIS-2024000 and a Google Faculty Research Award.

\bibliographystyle{plainnat}
% \bibliography{references}

\begin{thebibliography}{10}

\bibitem[He et~al.(2016)]{he2016deep}
Kaiming He, Xiangyu Zhang, Shaoqing Ren, and Jian Sun.
\newblock Deep residual learning for image recognition.
\newblock In \textit{CVPR}, pages 770--778, 2016.

\bibitem[Liu et~al.(2019)]{liu2019darts}
Hanxiao Liu, Karen Simonyan, and Yiming Yang.
\newblock {DARTS}: Differentiable architecture search.
\newblock In \textit{ICLR}, 2019.

\bibitem[Pham et~al.(2018)]{pham2018efficient}
Hieu Pham, Melody Guan, Barret Zoph, Quoc Le, and Jeff Dean.
\newblock Efficient neural architecture search via parameter sharing.
\newblock In \textit{ICML}, pages 4095--4104, 2018.

\bibitem[Real et~al.(2019)]{real2019regularized}
Esteban Real, Alok Aggarwal, Yanping Huang, and Quoc~V Le.
\newblock Regularized evolution for image classifier architecture search.
\newblock In \textit{AAAI}, volume~33, pages 4780--4789, 2019.

\bibitem[Tan et~al.(2019)]{tan2019mnasnet}
Mingxing Tan, Bo~Chen, Ruoming Pang, Vijay Vasudevan, Mark Sandler, Andrew Howard, and Quoc~V Le.
\newblock {MnasNet}: Platform-aware neural architecture search for mobile.
\newblock In \textit{CVPR}, pages 2820--2828, 2019.

\bibitem[Vaswani et~al.(2017)]{vaswani2017attention}
Ashish Vaswani, Noam Shazeer, Niki Parmar, Jakob Uszkoreit, Llion Jones, Aidan~N Gomez, Lukasz Kaiser, and Illia Polosukhin.
\newblock Attention is all you need.
\newblock In \textit{NeurIPS}, pages 5998--6008, 2017.

\bibitem[Wu et~al.(2019)]{wu2019fbnet}
Bichen Wu, Xiaoliang Dai, Peizhao Zhang, Yanghan Wang, Fei Sun, Yiming Wu, Yuandong Tian, Peter Vajda, Yangqing Jia, and Kurt Keutzer.
\newblock {FBNet}: Hardware-aware efficient convnet design via differentiable neural architecture search.
\newblock In \textit{CVPR}, pages 10734--10742, 2019.

\bibitem[Zoph and Le(2017)]{zoph2017neural}
Barret Zoph and Quoc~V Le.
\newblock Neural architecture search with reinforcement learning.
\newblock In \textit{ICLR}, 2017.

\end{thebibliography}

\end{document}
